\documentclass[11pt,a4paper]{article}

% --- Packages ---
\usepackage[margin=.75in]{geometry}
\usepackage[utf8]{inputenc}
\usepackage[T1]{fontenc}
\usepackage{lmodern} % Fix for font not found error
\usepackage{titlesec}
\usepackage{enumitem}
\usepackage{hyperref}
\usepackage{xcolor}
\usepackage{fancyhdr}
\usepackage{comment}
\usepackage{cleveref}
\usepackage[bottom]{footmisc}
\usepackage{graphicx}
\usepackage{longtable}
\usepackage{booktabs}
\usepackage{array}
\usepackage{calc}
\usepackage{etoolbox}

% Standard Pandoc setup
\providecommand{\tightlist}{%
  \setlength{\itemsep}{0pt}\setlength{\parskip}{0pt}}

\crefformat{section}{\S#2#1#3}
\crefformat{subsection}{\S#2#1#3}
\crefformat{subsubsection}{\S#2#1#3}
\crefformat{figure}{#2Figure~#1#3}
\crefformat{footnote}{#2\footnotemark[#1]#3}

\linespread{0.9} % Riduce lo spazio tra le linee

% --- Colors ---
\definecolor{primary}{RGB}{0, 0, 0}
\definecolor{links}{HTML}{0000EE}
\definecolor{blue}{HTML}{26428B}

\hypersetup{
    colorlinks=true,
    linkcolor=links,
    urlcolor=links,
    citecolor=links,
}

% --- Section Formatting ---
\titleformat{\section}
  {\normalfont\Large\bfseries\color{blue}}{\thesection}{1em}{}
\titleformat{\subsection}
{\normalfont\large\bfseries\color{blue}}{\thesubsection}{1em}{}
\titleformat{\subsubsection}
{\normalfont\bfseries\color{blue}}{\thesubsubsection}{1em}{}

% --- Header and Footer ---
\pagestyle{fancy}
\fancyhead[L]{\small Giuseppe Capasso}
\fancyhead[R]{\small virtual-machine}
\fancyfoot[C]{\thepage}
\renewcommand{\headrulewidth}{0.4pt}

% --- Custom Title Command ---
\newcommand{\proposaltitle}[1]{
    \begin{center}
        \vspace*{0.1cm}
        {\LARGE \bfseries \color{blue} #1} \\[1.5ex]
        {\large \textbf{Giuseppe Capasso}} \\[1ex]
                \small{
            \vspace{0.1cm}
            \textbf{Materia:} virtual-machine\\
        }
            \end{center}
    \vspace{0.3cm}
    \hrule height 0.5pt
    \vspace{0.5cm}
}

\begin{document}

\proposaltitle{Multi VM setup}

\section{Multi VM setup}\label{multi-vm-setup}

\subsection{Description}\label{description}

In this activity, you will learn how to make two different virtual
machines communicate.

You will create a dedicated virtual network on your hypervisor and
attach both virtual machines to it.

\subsection{Virtual Box}\label{virtual-box}

If you have Virtual Box as your hypervisor, you will need to go into
\texttt{Files\ -\textgreater{}\ Tools\ -\textgreater{}\ Network\ tools}
and create a new host network by clicking on the \texttt{New} button.

Network Settings

Once you click on \texttt{New}, a new network will be created with a
dedicated IP LAN and a DHCP. Now you can use this network to make all
your virtual machines to communicate.

\subsection{Attaching Host to the created
VM}\label{attaching-host-to-the-created-vm}

\paragraph{Metasploitable}\label{metasploitable}

For Metasploitable2 is simple: just go into settings and choose from the
current NIC (it's not needed to enable a new one).

IPv4 settings

Reboot the virtual machine and now it should have an IP from the virtual
network you created.

\subsubsection{Kali Linux}\label{kali-linux}

For example, if you want to attach Kali Linux to the new Network, go
into VM settings, enable an extra adapter and choose the
\texttt{vboxnet0}.

Kali Network Settings

When you enable an extra NIC, you will likely have to configure it (by
default Kali configures only the first NIC). Start the virtual machine
and enable the connection like in the pictures below.

First, go to \texttt{Advanced\ Network\ Settings} and create a new
Profile.

Advanced Network Settings

Choose Ethernet when asked.
\includegraphics[width=5.83333in,height=3.03999in]{media/rId35.png}

Select \texttt{eth1} for the device.
\includegraphics[width=5.83333in,height=3.03999in]{media/rId38.png}

Make sure, DHCP is enabled.
\includegraphics[width=5.83333in,height=3.03999in]{media/rId41.png}

\subsection{Result}\label{result}

Now you have two virtual machines in the same virtual network and you
can make practice on your own!

Here it is the result of nmap from Kali Linux in the virtual network
created.

Nmap result

\end{document}
